\documentclass[10pt,a4paper]{amsart}
\usepackage[margin=19mm]{geometry}
\usepackage{amsmath,amssymb,mathtools}
\usepackage[scr=rsfs,cal=euler]{mathalfa}
\usepackage{xfrac}
\usepackage[unicode,hidelinks,bookmarks=false]{hyperref}
\newcommand{\ie}{i.e.\ }
\newcommand{\cf}{cf.\ }
\newcommand{\resp}{respectively\ }
\newcommand{\Subsection}{\subsection}
\providecommand{\nobreakdash}{\nobreak\hbox{-}}
\setlength{\emergencystretch}{2em}
\multlinegap0pt
\usepackage{xcolor}
\usepackage{amssymb}
\usepackage[shortlabels]{enumitem}
\usepackage{mathtools}
\newtheorem{lem}{Lemma}
\newtheorem{thm}{Theorem}
\newcommand{\R}{\mathbb{R}}
\newcommand{\X}{\mathcal{X}}
\newcommand{\lammin}{\lambda_{\textup{min}}}
\newcommand{\lin}{{\textup{lin}}}
\newcommand{\norm}[2]{
    \|#1\|_{#2}
}
\newcommand{\partv}[1]{
    \frac{\partial}{\partial #1}
}
\newcommand{\s}{\psi}
\title{Quantitative convergence of trained single layer neural networks to Gaussian processes}
\author{Eloy Mosig García}
\date{Source: arXiv:2509.24544v3 (CC BY 4.0)}
\begin{document}
\maketitle

\noindent\emph{Source and licence.} Quantitative convergence of trained single layer neural networks to Gaussian processes. Version arXiv:2509.24544v3 (CC BY 4.0), distributed by arXiv under the Creative Commons Attribution 4.0 International licence, http://creativecommons.org/licenses/by/4.0/, as recorded in the arXiv licence metadata for this exact version and shown on the abstract page https://arxiv.org/abs/2509.24544v3. Full licence notice: \texttt{licenses/arxiv-2509.24544-CC-BY-4.0.txt}.\par\smallskip

\noindent\textit{Editorial scope.} Counted unit: the article's thm thm:stime\_polinomiali (source0.tex lines 1560--1583). The article's complete original proof of it is delivered with it (source0.tex lines 2601--2763, one \texttt{proof} environment of 163 lines, which the article places away from the statement). No mathematics is written, repaired or re-derived: the statement and the proof are the article's own lines, in the article's own notation, and no retained line is substituted or re-flowed.  Retained uncounted as the source-local context the proof needs: lines 1408--1423; lines 1596--1609.  Nothing was omitted from any retained span, so the package carries no omission record.  No retained line contains a citation, so the wrapper carries no bibliography and no bibliography entry.  No retained line contains a figure environment, an image-inclusion command or drawing code, so no image is delivered and the package declares no asset dependency.  Editorial formatting only, in addition: an AMS \texttt{amsart} preamble that restates the article's own declarations and macros used by the retained lines, namely the environments \texttt{lem}, \texttt{thm} on separate counters, together with the standard packages those lines need.  The retained environments print in this document as Lemma 1, Theorem 1 in order of appearance, whereas the article prints the same items with the numbering of its own full document.  The complete native source of the article as distributed in the arXiv e-print for this exact version is bundled unchanged as \texttt{sources/185953/source0.tex}.
\begin{lem}[Gradients $f$ and explicit formulae for $\tilde k$ and $k$]
    \label{lemma1}
    The following hold for each $x,x'\in \R^{n_0}$:
    \begin{align}
        \nabla_{\theta^{(0)}}f(x,\theta) &= \frac{1}{\sqrt{n_1n_0}}x^\top\Phi'(\frac{1}{\sqrt{n_0}}x\theta^{(0)})\theta^{(1)} \in \R^{n_0}, \\
        \nabla_{\theta^{(1)}}f(x,\theta) &= \frac{1}{\sqrt{n_1}}\Phi(\frac{1}{\sqrt{n_0}}x\theta^{(0)}) \in \R^{n_1}.
    \end{align}
    Moreover, $\tilde k(x,x')$ is a diagonal $n_1 \times n_1$ matrix with
    \begin{align}
        \tilde k_{ii}(x,x') = \frac{1}{n_0}\sum_{u=1}^{n_0}x_ux'_u,
    \end{align}
    and $k(x,x')$ is a real function given by:
    \begin{align}
        k(x,x') &=\frac{1}{n_1n_0}\sum_{u=1}^{n_0} x_ux'_u \sum_{v=1}^{n_1}\Phi'(h_v(x))\Phi'(h_v(x'))(\theta^{(1)}_v)^2 + \frac{1}{n_1}\sum_{v=1}^{n_1}\Phi(h_v(x))\Phi(h_v(x')).
    \end{align}
\end{lem}

    \begin{lem}[Inequalities for $\theta^{(i)}_t$]
        \label{lem:gronwallbeta}
        Fix $u$ and $v$ with $1 \le u \le n_0$ and $1 \le v \le n_1$ and put $\X_u = ((x_i)_u)_{i=1}^n\in \R^n$.
        Let $\lammin$ and $\lammin^0$ be the smallest eigenvalues of $k_t$ and $k_0$, respectively.
        Then the following inequalities hold:
        \begin{align}
            (\theta^{(1)}_v)_t &\le (\theta_v^{(1)})_0 + \frac{\norm{\Phi}{\infty}\norm{f_0-y}{}}{\sqrt{n_1}}I_t(\lammin),\\
            (\theta^{(0)}_{uv})_t &\le (\theta_{uv}^{(0)})_0 + \frac{\norm{\Phi}{\infty}\norm{\Phi'}{\infty}\norm{f_0-y}{}^2\norm{\X_u}{}}{2n_1\sqrt{n_0}}I_t(\lammin)^2\\
        &\quad + \frac{\norm{\Phi'}{\infty}\norm{f_0-y}{}\norm{\X_u}{}}{\sqrt{n_1n_0}}(\theta^{(1)}_v)_0{}I_t(\lammin),\\
        (\overline \theta^{(1)}_v)_t &\le (\theta_v^{(1)})_0 + \frac{\norm{\Phi}{\infty}\norm{f_0-y}{}}{\sqrt{n_1}}I_t(\lammin^0),\\
        (\overline \theta^{(0)}_{uv})_t &\le (\theta_{uv}^{(0)})_0 + \frac{\norm{\Phi}{\infty}\norm{\Phi'}{\infty}\norm{f_0-y}{}^2\norm{\X_u}{}}{2n_1\sqrt{n_0}}I_t(\lammin^0)^2 \\
    &\quad + \frac{\norm{\Phi'}{\infty}\norm{f_0-y}{}\norm{\X_u}{}}{\sqrt{n_1n_0}}(\theta^{(1)}_v)_0{}I_t(\lammin^0).
\end{align}
    \end{lem} 

\begin{thm}
    \label{thm:stime_polinomiali}
    Assume that $\Phi$ and $\Phi'$ are bounded.
    Let $L(\X)$ be the Lipschitz constant of $\nabla_\theta f_0$, seen as a function of $\theta$,
    and let $\s(\theta_0) = \norm{f_0-y}{}$.
    Then, for each $t > 0$ there exist positive constants $A_1, \ldots, A_5,B_1, \ldots, B_9,C_1, \ldots, C_8$ and $C_9$ not depending on $n_1, n_0$ nor $t$ such that, $\mathbb P$-almost everywhere:
    \begin{align}
        \norm{y_t- f_t}{}^2 &\le \frac{A_0}{n_0}\norm{\theta^{(0)}_0\theta^{(1)}_0}{}^2 +\frac{A_1t^2}{n_0n_1}\norm{\theta^{(0)}_0}{}^2\s(\theta_0)^2 + \frac{A_2\norm{\theta^{(1)}_0}{}^2t^4}{n_1^2n_0}\s(\theta_0)^4 \\
        &\quad + \frac{A_3t^6}{n_1^2n_0}\s(\theta_0)^6 + \frac{A_4t^2}{n_1n_0}\norm{\theta^{(1)}_0}{}^4\s(\theta_0)^2 + \frac{A_5t^4\norm{\theta^{(1)}_0}{}^2}{n_1^2n_0}\s(\theta_0)^4,
    \end{align}
    \begin{align}
        \norm{f^\lin-\overline y_t}{}^2 &\le \frac{B_0}{n_1n_0}\norm{\theta_0^{(0)}\theta_0^{(1)} }{}^2 + \frac{B_1}{n_1^2n_0^2}\norm{\theta_0^{(0)}\theta_0^{(1)} }{}^2\s(\theta_0)^4t^4\\
        &\quad +\frac{B_2}{n_1^2n_0^2}\norm{\theta_0^{(0)}}{}^2\norm{\theta_0^{(1)} }{}^4\s(\theta_0)^2t^2+\frac{B_3}{n_1n_0^2}\norm{\theta_0^{(0)}\theta_0^{(1)} }{}^2\s(\theta_0)^2 t^2\\
        &\quad +\frac{B_4}{n_1^2n_0}\norm{\theta_0^{(1)} }{}^2\s(\theta_0)^4t^4+\frac{B_5}{n_1^2n_0}\norm{\theta_0^{(1)} }{}^4\s(\theta_0)^2t^2 \\
        &\quad +\frac{B_6}{n_1n_0}\norm{\theta_0^{(1)} }{}^2\s(\theta_0)^2t^2+\frac{B_7}{n_1^2n_0}\norm{\theta_0^{(0)} }{}^2\s(\theta_0)^4t^4 \\
        &\quad +\frac{B_8}{n_1^2n_0}\norm{\theta_0^{(0)} }{}^2\norm{\theta_0^{(1)} }{}^2\s(\theta_0)^2t^2+\frac{B_9}{n_1n_0}\norm{\theta_0^{(0)} }{}^2\s(\theta_0)^2t^2,
    \end{align}
    \begin{align}
        \norm{f_t- f^\lin}{}^2 &\le \frac{L(\X)^2}{n_1^2}\left(\frac{C_1 \s(\theta_0)^4\norm{\theta^{(1)}_0}{}^2t^2}{n_1n_0^2}+\frac{C_2 \s(\theta_0)^6t^8}{n_1^2n_0^2}+\frac{C_3 \s(\theta_0)^4t^6}{n_1n_0}\right. \\
    &\left.\quad +\frac{C_4 \norm{\theta^{(1)}_0}{}^4t^4}{n_0^2}+\frac{C_5 \s(\theta_0)^2\norm{\theta^{(1)}_0}{}^2t^6}{n_1n_0^2}+\frac{C_6 \norm{\theta^{(1)}_0}{}^2t^6}{n_0}\right.\\
    &\left.\quad +\frac{C_7 \s(\theta_0)^2\norm{\theta^{(1)}_0}{}^2t^4}{n_0}+\frac{C_8 \s(\theta_0)^4t^6}{n_1n_0}+C_9\s(\theta_0)^2t^4\right),
\end{align}
where $\theta^{(0)}\theta^{(1)} \in \R^{n_0}$ denotes the usual product of the matrices $\theta^{(0)}$ and $\theta^{(1)}$.
\end{thm}

\begin{proof}[Proof of Theorem \ref{thm:stime_polinomiali}]
    We prove the three inequalities separately.
    Let $\lammin$ denote the smallest eigenvalue of $k_t$.
    \begin{itemize}
        \item By Lemma \ref{lemma1} and Cauchy-Schwarz's inequality,
        \begin{align}
            \label{eqaux_3triangle}
            \norm{y_t- f_t}{} &\le \frac{1}{\sqrt{n_1n_0}} \textup{Lip}\Phi \norm{x-\X}{} \norm{\theta_t^{(0)} \theta_t^{(1)} }{}.
        \end{align}
        Recall that $I_t(\lammin) \le t$.
        % and $I_t(I_t((\lammin))) \le 1$.
        Then the norm $\norm{\theta_t^{(0)} \theta_t^{(1)} }{}^2$ can be estimated with the aid of Lemma \ref{lem:gronwallbeta}:
        \begin{align}
            \norm{\theta_t^{(0)} \theta_t^{(1)}}{}^2 &= \sum_{u=1}^{n_0}\left(\sum_{v=1}^{n_1}(\theta^{(0)}_{uv})_t(\theta^{(1)}_v)_t\right)^2 \\
            &\le \sum_{u=1}^{n_0}\left(\sum_{v=1}^{n_1}(\theta_v^{(1)})_0(\theta_{uv}^{(0)})_0 + \frac{a_1(\theta_{uv}^{(0)})_0}{\sqrt{n_1}}\s(\theta_0)t + \frac{a_0(\theta_v^{(1)})_0}{n_1\sqrt{n_0}}\s(\theta_0)^2t^2 \right. \\
            &\quad \left. + \frac{a_0a_1}{n_1^\frac{3}{2}\sqrt{n_0}}\s(\theta_0)^3t^3 + \frac{a'_0(\theta_v^{(1)})^2_0}{\sqrt{n_1n_0}}\s(\theta_0)t + \frac{a'_0a_1(\theta^{(1)}_v)_0}{n_1\sqrt{n_0}}\s(\theta_0)^2t^2\right)^2\\
            &\le \sum_{u,v} n_1(\theta_v^{(1)})_0^2(\theta_{uv}^{(0)})_0^2 +a_1^2(\theta_{uv}^{(0)})^2_0\s(\theta_0)^2t^2 + \frac{a_0^2(\theta_v^{(1)})^2_0}{n_1n_0}\s(\theta_0)^4t^4 \\
            &\quad + \frac{a_0^2a_1^2}{n_1^2n_0}\s(\theta_0)^6t^6 + \frac{{a'}^2_0(\theta_v^{(1)})^4_0}{n_0}\s(\theta_0)^2t^2 + \frac{{a'}^2_0a_1^2(\theta^{(1)}_v)^2_0}{n_1n_0}\s(\theta_0)^4t^4 \\
            &\le n_1\norm{\theta^{(0)}_0\theta^{(1)}_0}{}^2 +a_1^2\norm{\theta^{(0)}_0}{}^2\s(\theta_0)^2t^2 + \frac{a_0^2\norm{\theta^{(1)}_0}{}^2}{n_1}\s(\theta_0)^4t^4 \\
            &\quad + \frac{a_0^2a_1^2}{n_1}\s(\theta_0)^6t^6 + {a'}^2_0\norm{\theta^{(1)}_0}{}^4\s(\theta_0)^2t^2 + \frac{{a'}^2_0a_1^2\norm{\theta^{(1)}_0}{}^2}{n_1}\s(\theta_0)^4t^4.
        \end{align}
        with $a_0 = \frac{1}{2}\norm{\Phi}{\infty}\norm{\Phi'}{\infty}\norm{\X_u}{}$, $a'_0 = \norm{\Phi'}{\infty}\norm{\X_u}{}$ and $a_1 = \norm{\Phi}{\infty}$.
        %with $\tilde a_1 = a_1^2$, 
        %$\tilde a_2=a_0^2$, 
        %$\tilde a_3=a_0^2a_1^2$,
        %$\tilde a_4={a'}^2_0$ and $\tilde a_5 = {a'}^2_0a_1^2$ positive constants not depending on $n_1$ nor $t$.

        Hence,
        \begin{align}
            \label{eqaux_1stsum}
            \norm{y_t- f_t}{}^2 &\le \frac{(\textup{Lip}\Phi)^2\norm{x-\X}{}}{n_0n_1}\norm{\theta^{(0)}_t\theta^{(1)}_t}{}^2\\
            &\le \frac{A_0}{n_0}\norm{\theta^{(0)}_0\theta^{(1)}_0}{}^2 +\frac{A_1t^2}{n_0n_1}\norm{\theta^{(0)}_0}{}^2\s(\theta_0)^2 + \frac{A_2\norm{\theta^{(1)}_0}{}^2t^4}{n_1^2n_0}\s(\theta_0)^4 \\
            &\quad + \frac{A_3t^6}{n_1^2n_0}\s(\theta_0)^6 + \frac{A_4t^2}{n_1n_0}\norm{\theta^{(1)}_0}{}^4\s(\theta_0)^2 + \frac{A_5t^4\norm{\theta^{(1)}_0}{}^2}{n_1^2n_0}\s(\theta_0)^4.
        \end{align}
        with $A_0 = (\textup{Lip}\Phi)^2\norm{x-\X}{}^2$,
        $A_1 = a_1^2$, $A_2= a_0^2$, $A_3 = a_1^2a_0^2$, $A_4 = {a'}_0^2$ and $A_5 = {a'}_0^2a_1^2$.

\item We follow a similar strategy to prove the second inequality in the Theorem. 
Put $\overline w_t = \overline \theta_t - \theta_0$.
By the triangle inequality and Cauchy-Schwarz we decompose:

\begin{align}
    \label{eqaux_flin}
        \norm{f^\lin-\overline y_t}{}^2 \le 2\norm{f_0 - y_0}{}^2 + 2\norm{\nabla_\theta f_0 - \nabla_\theta y_0}{}^2\norm{\overline w_t}{}^2.
    \end{align}
    The first summand in \eqref{eqaux_flin} is bounded exactly as the first summand in \eqref{eqaux_3triangle} by setting $t=0$:
    \begin{align}
        \norm{f_0 - y_0}{}^2 &\le  \frac{(\textup{Lip}\Phi)^2 \norm{x-\X}{}^2}{n_1n_0} \norm{\theta_0^{(0)}\theta_0^{(1)} }{}^2 \le \frac{A_0}{n_1n_0}\norm{\theta_0^{(0)}\theta_0^{(1)} }{}^2.
    \end{align}

    As for the second summand in \eqref{eqaux_flin}, we decompose by Lemma \ref{lemma1}.
    Factoring out $\max_i \{(\theta^{(1)}_0)_i\} = \norm{\theta^{(1)}_0}{\infty}$ permits us to write:
    \begin{align}
        \norm{\nabla_\theta f_0 - \nabla_\theta y_0}{}^2 &\le \norm{\partv{\theta^{(0)}}(f_0 - y_0)}{}^2+\norm{\partv{\theta^{(1)}}(f_0 - y_0)}{}^2 \\
        &\le \frac{1}{n_1n_0} \norm{(\X^\top\Phi'(\frac{1}{\sqrt{n_0}}\X\theta_0^{(0)})-x^\top\Phi'(\frac{1}{\sqrt{n_0}}x\theta_0^{(0)}))\theta_0^{(1)}}{}^2\\
        &\quad + \frac{1}{n_1}\norm{\Phi(\frac{1}{\sqrt{n_0}}\X\theta_0^{(0)})-\Phi(\frac{1}{\sqrt{n_0}}x\theta_0^{(0)})}{}^2 \\
        &\le \frac{2}{n_1n_0} \norm{(\X^\top\Phi'(\frac{1}{\sqrt{n_0}}\X\theta_0^{(0)})-\X^\top\Phi'(\frac{1}{\sqrt{n_0}}x\theta_0^{(0)}))\theta_0^{(1)}}{}^2\\
        &\quad +\frac{2}{n_1n_0} \norm{(\X^\top\Phi'(\frac{1}{\sqrt{n_0}}x\theta_0^{(0)})-x^\top\Phi'(\frac{1}{\sqrt{n_0}}x\theta_0^{(0)}))\theta_0^{(1)}}{}^2\\
        &\quad + \frac{A_0}{n_1n_0}\norm{\theta_0^{(0)}}{}^2 \\
        &\le \frac{2}{n_1n_0^2}(\textup{Lip}\Phi')^2\norm{\X}{}^2 \norm{x-\X}{}^2\norm{\theta_0^{(0)}\theta_0^{(1)}}{}^2\\
        &\quad +\frac{2}{n_1n_0} \norm{\Phi'}{\infty}^2\norm{x-\X}{}^2\norm{\theta_0^{(1)}}{}^2 + \frac{A_0}{n_1n_0}\norm{\theta_0^{(0)}}{}^2.
    \end{align}

    Moreover we can bound the norm of $\overline w_t$ with Lemma \ref{lem:gronwallbeta}:
    \begin{align}
        \label{eqaux_flinaux2}
        \norm{\overline w_t}{}^2 &\le \norm{\overline \theta_t^{(0)}-\theta_0^{(0)}}{}^2 + \norm{\overline \theta_t^{(1)}-\theta_0^{(1)}}{}^2 \\
        &\le \sum_{u,v}\frac{2a_0^2}{n_1^2n_0}\s(\theta_0)^4t^4 + \frac{2{a'}^2_0(\theta^{(1)}_v)^2_0}{n_1n_0}\s(\theta_0)^2t^2+ \sum_{v} \frac{a_1^2}{n_1}\s(\theta_0)^2t^2\\
        &\le \frac{2a_0^2}{n_1}\s(\theta_0)^4t^4 + \frac{2{a'}^2_0\norm{\theta^{(1)}}{}^2}{n_1}\s(\theta_0)^2t^2+ a_1^2\s(\theta_0)^2t^2.
    \end{align}

    Hence, \eqref{eqaux_flin} can be written as:
    \begin{align}
        \label{eqaux_3rdsum} \norm{f^\lin-\overline y_t}{}^2 &\le \frac{B_0}{n_1n_0}\norm{\theta_0^{(0)}\theta_0^{(1)} }{}^2 + \frac{B_1}{n_1^2n_0^2}\norm{\theta_0^{(0)}\theta_0^{(1)} }{}^2\s(\theta_0)^4t^4\\
        &\quad +\frac{B_2}{n_1^2n_0^2}\norm{\theta_0^{(0)}}{}^2\norm{\theta_0^{(1)} }{}^4\s(\theta_0)^2t^2+\frac{B_3}{n_1n_0^2}\norm{\theta_0^{(0)}\theta_0^{(1)} }{}^2\s(\theta_0)^2 t^2\\
        &\quad +\frac{B_4}{n_1^2n_0}\norm{\theta_0^{(1)} }{}^2\s(\theta_0)^4t^4+\frac{B_5}{n_1^2n_0}\norm{\theta_0^{(1)} }{}^4\s(\theta_0)^2t^2 \\
        &\quad +\frac{B_6}{n_1n_0}\norm{\theta_0^{(1)} }{}^2\s(\theta_0)^2t^2+\frac{B_7}{n_1^2n_0}\norm{\theta_0^{(0)} }{}^2\s(\theta_0)^4t^4 \\
        &\quad +\frac{B_8}{n_1^2n_0}\norm{\theta_0^{(0)} }{}^2\norm{\theta_0^{(1)} }{}^2\s(\theta_0)^2t^2+\frac{B_9}{n_1n_0}\norm{\theta_0^{(0)} }{}^2\s(\theta_0)^2t^2,
    \end{align}
    where the constants in the last inequality are, explicitly, $B_0=2A_0$,
    $B_1=8(\textup{Lip}\Phi'\norm{\X}{}\norm{x-\X}{}a_0)^2$,
    $B_2=8(\textup{Lip}\Phi'\norm{\X}{}\norm{x-\X}{}a'_0)^2$,
    $B_3=4(\textup{Lip}\Phi'\norm{\X}{}\norm{x-\X}{}a_1)^2$,
    $B_4=8(\norm{\Phi'}{\infty}\norm{x-\X}{}a_0)^2$,
    $B_5=8(\norm{\Phi'}{\infty}\norm{x-\X}{}a'_0)^2$,
    $B_6=4(\norm{\Phi'}{\infty}\norm{x-\X}{}a_1)^2$,
    $B_7=4A_0a_0^2$,
    $B_8=4A_0{a'}_0^2$, and $B_9 = 2A_0a_1^2$.

    \item It remains to estimate the last inequality.
        Consider $\Delta(t) = \norm{f_t - f_t^\lin}{}$
        Then by gradient flow equations and Cauchy-Schwarz,
        \begin{align}
            \partv{t}(\Delta(t)^2) & = \langle k_t(f_t-y) -k_0(f_t^\lin-y), f_t - f_t^\lin \rangle \\
            &= \sum_{i=1}^n (k_t(x_i,\X)(f_t-y) -k_0(x_i,\X)(f_t^\lin-y))(f_t(x_i) - f_t^\lin(x_i))\\
            &= \sum_{i=1}^n (k_t(x_i,\X)-k_0(x_i,\X))(f_t-y)(f_t(x_i) - f_t^\lin(x_i)) \\
            &\quad -k_0(x_i,\X)(f_t - f_t^\lin)(f_t(x_i) - f_t^\lin(x_i))\\
            &=\norm{(k_t-k_0)(f_t-y)(f_t-f_t^\lin)^\top}{1} - \norm{k_0(f_t-f_t^\lin)(f_t-f_t^\lin)^\top}{1}
        \end{align}
        By equivalence of the $1$-norm and the euclidean norm for $v\in \R^d$ we have $\norm{v}{} \le \norm{v}{1} \le \sqrt{d}\norm{v}{}$.
        Then, by Cauchy-Schwarz's inequality,
        \begin{align}
            \frac{\partial}{\partial t} \Delta(t) &= n\norm{k_t-k_0}{}\norm{f_t-y}{} -\lammin^0\Delta(t) \\
            &\label{eqaux_delta} \le n\norm{k_t-k_0}{}\s(\theta_0)e^{-\lammin t} -\lammin^0\Delta(t)
        \end{align}

        Let us bound the norm of $k_t-k_0$:
        \begin{align}
            \norm{k_t-k_0}{}&= \norm{\nabla_\theta f_t(\X)\nabla_\theta f_t(\X)^\top-\nabla_\theta f_0(\X)\nabla_\theta f_0(\X)^\top }{}\\
            &\le \norm{\nabla_\theta f_t(\X)+\nabla_\theta f_0(\X)}{}L(\X)\norm{\theta_t-\theta_0}{}\\
        \end{align}

        From Lemmas \ref{lemma1} and \ref{lem:gronwallbeta}, we have:
        \begin{align}
            \norm{\nabla_\theta f_t(\X)}{}^2 &=\norm{\nabla_{\theta^{(0)}}f_t(\X)}{}^2+\norm{\nabla_{\theta^{(1)}} f_t(\X)}{}^2\\
            &\le \frac{\norm{\X}{}^2\norm{\Phi'}{\infty}^2\norm{\theta^{(1)}_t}{}^2}{n_1n_0} + \frac{\norm{\Phi}{\infty}^2}{n_1}\\
            &\label{eqaux_delta1}\le \frac{\norm{\X}{}^2\norm{\Phi'}{\infty}^2}{n_1n_0}\left(\norm{\theta^{(1)}_0}{} + \frac{\norm{\Phi}{\infty}\s(\theta^0)}{\sqrt{n_1}}t\right)^2 + \frac{\norm{\Phi}{\infty}^2}{n_1}.
        \end{align}

        Analogously,
        \begin{align}
            \norm{\nabla_\theta f_0(\X)}{}^2 &=\norm{\nabla_{\theta^{(0)}}f_0(\X)}{}^2+\norm{\nabla_{\theta^{(1)}} f_0(\X)}{}^2\\
            &\label{eqaux_delta2}\le \frac{\norm{\X}{}^2\norm{\Phi'}{\infty}^2\norm{\theta^{(1)}_0}{}^2}{n_1n_0} + \frac{\norm{\Phi}{\infty}^2}{n_1}.
        \end{align}

        Moreover, again by Lemma \ref{lem:gronwallbeta},
        \begin{align}
            \norm{\theta_t-\theta_0}{}^2 &= \norm{\theta_t^{(0)}-\theta_0^{(0)}}{}^2 + \norm{\theta_t^{(1)}-\theta_0^{(1)}}{}^2 \\
            &\label{eqaux_delta3}\le \frac{2a_0^2}{n_1^2n_0}\s(\theta_0)^4t^4 + \frac{2{a'_0}^2}{n_1n_0}\norm{\theta^{(1)}_0}{}^2t^2 +\frac{a_1^2\s(\theta^0)^2}{n_1}t^2 \\
        \end{align}

        Inequalities \eqref{eqaux_delta1},\eqref{eqaux_delta2} and \eqref{eqaux_delta3} allow us to estimate:
        \begin{align}
            \label{eqaux_monster}
            \norm{k_t-k_0}{}&\le \frac{L(\X)}{n_1}\left(\frac{c_1 \s(\theta_0)^2\norm{\theta^{(1)}_0}{}}{\sqrt{n_1}n_0}+\frac{c_2 \s(\theta_0)^3t^3}{n_1n_0}+\frac{c_3 \s(\theta_0)^2t^2}{\sqrt{n_1n_0}} \right.\\
            &\left.\quad +\frac{c_4 \norm{\theta^{(1)}_0}{}^2t}{n_0}+\frac{c_5 \s(\theta_0)\norm{\theta^{(1)}_0}{}t^2}{\sqrt{n_1}n_0}+\frac{c_6 \norm{\theta^{(1)}_0}{}t}{\sqrt{n_0}}\right.\\
            &\left.\quad +\frac{c_7 \s(\theta_0)\norm{\theta^{(1)}_0}{}t}{\sqrt{n_0}}+\frac{c_8 \s(\theta_0)^2t^2}{\sqrt{n_1n_0}}+c_9\s(\theta_0)t\right),
        \end{align}
        with $c_1 = 2\sqrt{2}a_0\norm{\X}{}\norm{\Phi'}{\infty}$,
$c_2 = \sqrt{2}a_0\norm{\Phi}{\infty}$,
$c_3 = 2\sqrt{2}a_0\norm{\Phi}{\infty}$,
$c_4 = 2\sqrt{2}a'_0\norm{\X}{}\norm{\Phi'}{\infty}$,
$c_5 = \sqrt{2}a'_0\norm{\Phi}{\infty}$,
$c_6 = 2\sqrt{2}a'_0\norm{\Phi}{\infty}$,
$c_7 = 2a_1\norm{\X}{}\norm{\Phi'}{\infty}$,
$c_8 = a_1\norm{\Phi}{\infty}$ and $c_9 = 2a_1\norm{\Phi}{\infty}$.
Let $C(n_1,n_0,t,\theta_0)$ be the right hand side of \eqref{eqaux_monster}.
Then, the reight-hand side of \eqref{eqaux_delta} can be $\mathbb P$-almost surely bounded from above with:
\begin{align}
    \frac{\partial}{\partial t} \Delta(t) &\le n\norm{k_t-k_0}{}\s(\theta_0)e^{-\lammin t} -\lammin^0\Delta(t) \\
    &\le nC(n_1,n_0,t,\theta_0).
\end{align}
In the previous inequality we used that the event $\lammin=0$ has null measure.
Integrating, and using that $\Delta(0)=0$:
\begin{align}
    \Delta(t) &\le \frac{nL(\X)}{n_1}\left(\frac{c_1 \s(\theta_0)^2\norm{\theta^{(1)}_0}{}t}{\sqrt{n_1}n_0}+\frac{c_2 \s(\theta_0)^3t^4}{2n_1n_0}+\frac{c_3 \s(\theta_0)^2t^3}{\sqrt{n_1n_0}}\right. \\
    &\left.\quad +\frac{c_4 \norm{\theta^{(1)}_0}{}^2t^2}{2n_0}+\frac{c_5 \s(\theta_0)\norm{\theta^{(1)}_0}{}t^3}{3\sqrt{n_1}n_0}+\frac{c_6 \norm{\theta^{(1)}_0}{}t^3}{2\sqrt{n_0}}\right.\\
    &\left.\quad +\frac{c_7 \s(\theta_0)\norm{\theta^{(1)}_0}{}t^2}{2\sqrt{n_0}}+\frac{c_8 \s(\theta_0)^2t^3}{3\sqrt{n_1n_0}}+\frac{c_9\s(\theta_0)t^2}{2}\right)
\end{align}
Taking the square, applying the elementary inequality $(\sum^n_{i=1}a_i)^2\le n\sum^n_{i=1}a_i^2$, for $a_i\ge 0$,
and adjusting the constants yields the desired result.
    \end{itemize}
\end{proof}

\end{document}
